\documentclass[a4paper,11pt]{article}

\usepackage[T1]{fontenc}
\usepackage[utf8]{inputenc}
\usepackage{lmodern}
\usepackage[margin=25mm]{geometry}
\usepackage{microtype}
\usepackage{amsmath,amssymb}
\usepackage{graphicx}
\usepackage{booktabs}
\usepackage{siunitx}
\usepackage[backend=biber,style=numeric,sorting=none]{biblatex}
\addbibresource{references.bib}
\usepackage[hidelinks]{hyperref}

\graphicspath{{figures/selected/}}
\sisetup{detect-all}
\setlength{\parindent}{0pt}
\setlength{\parskip}{0.5\baselineskip}
\setlength{\emergencystretch}{2em}

% EDIT: Replace the student name, experiment title and experiment date.
\newcommand{\studentname}{[Student name]}
\newcommand{\experimenttitle}{Observing Single-Photon Interference}
\newcommand{\experimentdate}{[Experiment date]}
\title{\experimenttitle}
\author{\studentname\\[0.4em]
  \small Queen's University Belfast\\
  \small PHY2001 --- Quantum and Statistical Physics}
\date{\small Experiment date: \experimentdate}

\begin{document}
\maketitle

\begin{abstract}
  % ADD ABSTRACT: Summarise the aim, approach, principal measured findings
  % (with uncertainties where appropriate) and conclusion after writing the report.
  \noindent\textit{[Abstract to be written by the student.]}
\end{abstract}

% Aim for approximately five pages in the completed report. No forced page
% breaks or table of contents are needed for a report of this length.
\section{Introduction and Theoretical Background}
\label{sec:introduction}

% ADD SCIENTIFIC EXPLANATION: State the aim and introduce the relevant theory
% in your own words. Cite the sources used for background and equations.
% ADD EQUATIONS: Define symbols, assumptions and applicable conditions.

% NUMBERED EQUATION EXAMPLE (replace placeholders before uncommenting):
% \begin{equation}
%   \text{[Insert equation here]}
%   \label{eq:theory}
% \end{equation}
% Refer to Equation~\eqref{eq:theory} in the accompanying explanation.

% BIBLIOGRAPHIC CITATION EXAMPLE:
% Add a matching entry to references.bib, then cite it with \cite{source-key}.

\section{Experimental Method}
\label{sec:methods}

% ADD METHOD: Describe the apparatus, alignment, acquisition procedure and
% instrument settings needed to reproduce your measurements.
% ADD MEASUREMENTS: Identify measured quantities, units, instrument resolution
% and how measurement uncertainties were assessed.
% ADD FIGURE: Include a labelled apparatus diagram if it aids reproducibility.

\section{Results and Analysis}
\label{sec:results}

% ADD RESULTS: Present your own measurements and analysis only.
% ADD UNCERTAINTY CALCULATIONS: Explain propagation, fitting and the meaning
% of reported uncertainties. Give units and consistent significant figures.

\subsection{Laser Wavelength Determination}
\label{sec:wavelength}
% ADD MEASUREMENTS, EQUATIONS AND ANALYSIS: Insert your measurements, analysis
% method, fitted quantities, wavelength estimate and uncertainty calculation.

\subsection{Laser Attenuation and Mean Photon Occupation}
\label{sec:attenuation}
% ADD MEASUREMENTS AND EQUATIONS: Record attenuation settings and measured
% quantities; show your occupation calculation, assumptions and uncertainty.

\subsection{Photon-Counting Statistics and Fano-Factor Analysis}
\label{sec:statistics}
% ADD TABLES, GRAPHS AND ANALYSIS: Present count data, acquisition intervals,
% statistical calculations and your Fano-factor estimate with uncertainty.

\subsection{Low-Flux Double-Slit Interference}
\label{sec:interference}
% ADD GRAPHS AND ANALYSIS: Present your measured interference data with units,
% error bars, any fitted model and the corresponding uncertainty calculations.

% FIGURE EXAMPLE: Save the actual graph in figures/selected/ before enabling.
% \begin{figure}[htbp]
%   \centering
%   \includegraphics[width=0.8\linewidth]{your-graph.pdf}
%   \caption{[Describe the measured quantities, conditions and error bars.]}
%   \label{fig:your-graph}
% \end{figure}
% Refer to Figure~\ref{fig:your-graph} in the results text.

% BOOKTABS TABLE EXAMPLE: Replace placeholders with your own measurements.
% \begin{table}[htbp]
%   \centering
%   \caption{[Describe the measurements and uncertainty convention.]}
%   \label{tab:measurements}
%   \begin{tabular}{lll}
%     \toprule
%     Quantity & Value (unit) & Uncertainty (unit) \\
%     \midrule
%     [Quantity] & [Value] & [Uncertainty] \\
%     \bottomrule
%   \end{tabular}
% \end{table}
% Refer to Table~\ref{tab:measurements} in the results text.
% SI UNIT SYNTAX: \si{\metre} for units; use \SI{<value>}{\metre}
% only after replacing <value> with a measured numerical value.

\section{Discussion and Critical Evaluation}
\label{sec:discussion}

% ADD SCIENTIFIC EXPLANATION: Interpret your results against the stated aims
% and theory. Assess agreement using uncertainties and cite comparisons.
% ADD CRITICAL EVALUATION: Discuss the quality of fits, statistical evidence,
% assumptions and any discrepancies supported by your measurements.

\section{Experimental Limitations and Improvements}
\label{sec:limitations}

% ADD LIMITATIONS: Identify relevant random and systematic effects and assess
% their influence on the reported results using evidence from your experiment.
% ADD IMPROVEMENTS: Propose specific changes and explain how each would
% address an identified limitation or improve measurement uncertainty.

\section{Conclusion}
\label{sec:conclusion}

% ADD CONCLUSION: Answer the experimental aim concisely using your principal
% measured results and uncertainties. State what your evidence supports.


% The heading remains visible before any references have been added.
\printbibheading[heading=bibnumbered,title={References}]
\printbibliography[heading=none]

\end{document}
